\section{A cubical Whitney counterpart}

We conclude with the cubical counterpart on a single cell. The purpose of this
section is limited: it proves the one-rectangular-cell cochain analogue of the
box boundary readout theorem and then records the local Whitney transfer needed
to compare that finite-dimensional statement with the preparatory homotopy
estimate.  It does not claim a global cubical-complex homotopy, a multi-cell
patching theorem, or inter-cell trace compatibility. First, the oriented
coboundary on one top-dimensional cube gives the exact discrete analogue of the
boundary readout theorem on the unit cube. The Whitney construction is then
obtained by composing the continuous homotopy operator with the classical
Whitney map and the integration map. Together these give the one-cell discrete
versions of the boundary-determination and decomposition mechanisms; compare
\cite{Whitney1957,Dodziuk1976,Dupont1976,DupontBook1978,Albert2020,
ArnoldFalkWinther2010,ArnoldAwanou2014,KaczynskiMischaikowMrozek2004}. The
results in this section are one-cell statements, including
Theorem~\ref{thm:cubical-boundary-readout} and
Proposition~\ref{prop:discrete-stokes-poincare}. Extending them to a multi-cell
cubical complex requires compatibility of the local homotopies and traces across
adjacent cells and is intentionally outside the present paper.

The one-cell boundary has an additional feature absent from the continuous
trace problem: it is a finite weighted atomic space.  The calibration deficit
therefore determines not only a linear error bound but the exact largest trace
error at every prescribed norm deficit.  The maximizing sign pattern can
change only when two explicit affine functions cross, producing a finite
phase-transition profile.

Let $C^k(I^n;\RR)$ be the space of cubical $k$-cochains on the standard cell
decomposition of $I^n$, with coboundary $\cob$. The cochain complex used for a
single rectangular cell $R=\prod_j[0,L_j]$ below is the finite-dimensional
complex generated by that oriented top cell and its oriented boundary faces,
with the fixed coordinate orientations displayed in the theorem; no additional
subdivision or global cubical complex is part of the notation. For
$\theta\in C^k(I^n;\RR)$ we set
\[
\norm{\theta}_{\cell,\infty}:=
\max_{\sigma^{(k)}}\abs{\theta(\sigma)},
\]
where the maximum is taken over all oriented $k$-cells.


\begin{theorem}[Cubical boundary readout on one rectangular cell]
\label{thm:cubical-boundary-readout}
Let $k\ge 1$, let $L_1,\dots,L_k>0$, and set
\[
R:=\prod_{j=1}^k[0,L_j],
\qquad
A_j:=\prod_{i\ne j}L_i,
\qquad
\abs{R}:=\prod_{j=1}^kL_j.
\]
Set also
\[
P_1(R):=2\sum_{j=1}^kA_j,
\qquad
m_R:=\frac{1}{2\sum_{j=1}^kL_j^{-1}}.
\]
Let $\Theta_R$ be the top cochain with value $\abs{R}$ on the oriented
rectangular cell $R$. For $1\le j\le k$ and $\lambda\in\{0,L_j\}$, let
$F_{j,\lambda}:=\{x_j=\lambda\}$ be oriented by
$\dd x_1\wedge\cdots\wedge\widehat{\dd x_j}\wedge\cdots\wedge\dd x_k$, and set
\[
s_{j,L_j}:=(-1)^{j-1},
\qquad
s_{j,0}:=-(-1)^{j-1}.
\]
For a boundary $(k-1)$-cochain $\xi$, define the weighted face norm by
\[
\norm{\xi}_{\cell,L,\infty}:=
\max_{1\le j\le k,\ \lambda\in\{0,L_j\}}
\frac{\abs{\xi(F_{j,\lambda})}}{A_j}.
\]
Let
\[
\mathcal F_R:=\{(j,\lambda):1\le j\le k,\
\lambda\in\{0,L_j\}\},
\qquad
w_{j,\lambda}:=A_j,
\]
and, for $S\subset\mathcal F_R$, write
\[
w(S):=\sum_{(j,\lambda)\in S}w_{j,\lambda}.
\]
For $\delta\ge0$ define
\begin{equation}
\label{eq:cubical-sharp-profile}
\Phi_R(\delta):=
2\max_{\varnothing\ne S\subsetneq\mathcal F_R}
\min\left\{
\delta\bigl(P_1(R)-w(S)\bigr),
\bigl(2m_R+\delta\bigr)w(S)
\right\}.
\end{equation}
Define the canonical boundary cochain $\eta_{\cell,R}\in C^{k-1}(R;\RR)$ by
\[
\eta_{\cell,R}(F_{j,\lambda})
=
s_{j,\lambda}m_RA_j
\qquad
(1\le j\le k,\ \lambda\in\{0,L_j\}).
\]
Then $\cob\eta_{\cell,R}=\Theta_R$ and
$\norm{\eta_{\cell,R}}_{\cell,L,\infty}=m_R$. Moreover, if
$\eta\in C^{k-1}(R;\RR)$ satisfies $\cob\eta=\Theta_R$ and
$M:=\norm{\eta}_{\cell,L,\infty}$, then
\begin{equation}
\label{eq:cubical-deficit-identity}
\sum_{j=1}^k\sum_{\lambda\in\{0,L_j\}}
A_j\left(M-s_{j,\lambda}\frac{\eta(F_{j,\lambda})}{A_j}\right)
=
P_1(R)(M-m_R).
\end{equation}
If $\delta:=M-m_R$, then the exact sharp stability estimate is
\begin{equation}
\label{eq:cubical-boundary-stability}
\sum_{j=1}^k\sum_{\lambda\in\{0,L_j\}}
A_j\left\lvert
\frac{\eta(F_{j,\lambda})}{A_j}
-
\frac{\eta_{\cell,R}(F_{j,\lambda})}{A_j}
\right\rvert
\le
\Phi_R(\delta).
\end{equation}
For every $\delta\ge0$ there is a boundary cochain with
$\cob\eta=\Theta_R$, norm $m_R+\delta$, and equality in
\eqref{eq:cubical-boundary-stability}.  Thus $\Phi_R$ is the exact stability
profile, not only an upper bound.  In particular, if
\[
A_{\min}:=\min_{1\le j\le k}A_j,
\qquad
0\le\delta\le
\frac{2m_RA_{\min}}{P_1(R)-2A_{\min}},
\]
where the quotient is interpreted as $+\infty$ when its denominator is zero,
then
\begin{equation}
\label{eq:cubical-small-deficit-profile}
\Phi_R(\delta)=2\bigl(P_1(R)-A_{\min}\bigr)\delta.
\end{equation}
In particular, $M\ge m_R$, the anisotropic constant $m_R$ is sharp, and every
minimizing cochain with $\cob\eta=\Theta_R$ has the canonical weighted boundary
values
\[
\eta(F_{j,\lambda})=\eta_{\cell,R}(F_{j,\lambda})
\qquad
\text{for all }j,\lambda.
\]
For $L_1=\cdots=L_k=1$ this reduces to the unit-cube constant $1/(2k)$ and
\begin{equation}
\label{eq:cubical-unit-profile}
\Phi_{I^k}(\delta)
=2\max_{1\le r\le 2k-1}
\min\left\{\delta(2k-r),(k^{-1}+\delta)r\right\}.
\end{equation}
For $k\ge2$ and $0\le\delta\le1/(2k(k-1))$, the sharp bound is
$\Phi_{I^k}(\delta)=(4k-2)\delta$; for $k=1$ it is $2\delta$ for every
$\delta\ge0$. In particular, on the unit $(k+1)$-cube the normalized
top cochain of cell norm one has optimal primitive norm $1/(2(k+1))$.
\end{theorem}

\begin{proof}
The boundary of the oriented top rectangular cell is
\[
\partial R
=
\sum_{j=1}^k\sum_{\lambda\in\{0,L_j\}}
s_{j,\lambda}F_{j,\lambda}.
\]
Therefore
\[
\cob\eta_{\cell,R}(R)
=
\eta_{\cell,R}(\partial R)
=
\sum_{j=1}^k\sum_{\lambda\in\{0,L_j\}}
s_{j,\lambda}^2m_RA_j
=2m_R\sum_{j=1}^kA_j
=\abs{R},
\]
because $A_j=\abs{R}/L_j$ and $2m_R\sum_jL_j^{-1}=1$. Thus
$\cob\eta_{\cell,R}=\Theta_R$. Its weighted face norm is plainly $m_R$.

Now let $\cob\eta=\Theta_R$. Since $\eta(\partial R)=\abs{R}$, we have
\[
\sum_{j=1}^k\sum_{\lambda\in\{0,L_j\}}
s_{j,\lambda}\eta(F_{j,\lambda})=\abs{R}.
\]
Subtracting this identity from $P_1(R)M=2\sum_jA_jM$ gives
\eqref{eq:cubical-deficit-identity}. Each summand on the left of
\eqref{eq:cubical-deficit-identity} is nonnegative because
$s_{j,\lambda}\eta(F_{j,\lambda})/A_j\le
\abs{\eta(F_{j,\lambda})}/A_j\le M$. This is precisely the finite-dimensional
linear-programming dual certificate for the weighted norm lower bound: with
$z_{j,\lambda}:=\eta(F_{j,\lambda})/A_j$, the primal problem is to minimize
$\max_{j,\lambda}|z_{j,\lambda}|$ subject to
$\sum_{j,\lambda}A_js_{j,\lambda}z_{j,\lambda}=\abs{R}$, while the dual functional
$z\mapsto P_1(R)^{-1}\sum_{j,\lambda}A_js_{j,\lambda}z_{j,\lambda}$ has dual norm one and value
$\abs{R}/P_1(R)=m_R$ on the affine constraint. The displayed deficit identity
is the associated complementary-slackness identity. Hence $M\ge m_R$.

Set $\delta:=M-m_R$ and, for $a=(j,\lambda)\in\mathcal F_R$, put
\[
y_a:=s_{j,\lambda}\frac{\eta(F_{j,\lambda})}{A_j},
\qquad
q_a:=y_a-m_R,
\qquad
w_a:=A_j.
\]
The coboundary constraint and $m_RP_1(R)=\abs{R}$ give
\begin{equation}
\label{eq:cubical-error-zero-mean}
\sum_{a\in\mathcal F_R}w_aq_a=0.
\end{equation}
The norm constraint $|y_a|\le M=m_R+\delta$ is equivalent to
\begin{equation}
\label{eq:cubical-error-caps}
-2m_R-\delta\le q_a\le\delta
\qquad(a\in\mathcal F_R).
\end{equation}
Moreover, the left-hand side of
\eqref{eq:cubical-boundary-stability} is
$\sum_aw_a|q_a|$.

Suppose first that $\delta>0$.  The vector $q$ is not identically zero,
because otherwise $M=m_R$, and hence its positive and negative supports are
both nonempty by \eqref{eq:cubical-error-zero-mean}. Let
\[
S:=\{a\in\mathcal F_R:q_a<0\},
\qquad
T:=\sum_{q_a>0}w_aq_a=-\sum_{q_a<0}w_aq_a.
\]
The two caps in \eqref{eq:cubical-error-caps} imply
\[
T\le\delta\bigl(P_1(R)-w(S)\bigr),
\qquad
T\le(2m_R+\delta)w(S).
\]
Since the weighted mean is zero,
\[
\sum_{a\in\mathcal F_R}w_a|q_a|=2T
\le
2\min\left\{
\delta\bigl(P_1(R)-w(S)\bigr),
(2m_R+\delta)w(S)
\right\}
\le\Phi_R(\delta).
\]
This proves \eqref{eq:cubical-boundary-stability}; the case $\delta=0$
follows directly from the same zero-mean identity and the upper cap $q_a\le0$.

It remains to prove that the profile is attained. Fix a set
$\varnothing\ne S\subsetneq\mathcal F_R$ attaining the maximum in
\eqref{eq:cubical-sharp-profile}, write $W:=w(S)$, and set
\[
T:=\min\{\delta(P_1(R)-W),(2m_R+\delta)W\}.
\]
Define
\begin{equation}
\label{eq:cubical-profile-extremizer}
q_a:=
\begin{cases}
-T/W,&a\in S,\\
T/(P_1(R)-W),&a\notin S.
\end{cases}
\end{equation}
Then \eqref{eq:cubical-error-zero-mean} holds, and the definition of $T$
gives the caps \eqref{eq:cubical-error-caps}.  Define a boundary cochain by
\[
\eta(F_{j,\lambda})
:=s_{j,\lambda}A_j(m_R+q_{j,\lambda}).
\]
The zero-mean identity shows that $\cob\eta=\Theta_R$. If the first term in
the definition of $T$ is minimal, every $q_a$ outside $S$ equals $\delta$; if
the second is minimal, every $q_a$ in $S$ equals $-2m_R-\delta$. Thus in
either case $\norm{\eta}_{\cell,L,\infty}=m_R+\delta$. Finally,
\[
\sum_aw_a|q_a|=2T=\Phi_R(\delta),
\]
so equality holds in \eqref{eq:cubical-boundary-stability}.

Every nonempty subset of $\mathcal F_R$ has weight at least $A_{\min}$.
Consequently each term in \eqref{eq:cubical-sharp-profile} is at most
$2\delta(P_1(R)-A_{\min})$. If
\[
\delta(P_1(R)-A_{\min})
\le(2m_R+\delta)A_{\min},
\]
which is precisely the range in
\eqref{eq:cubical-small-deficit-profile}, a singleton face of weight
$A_{\min}$ attains this upper bound. This proves
\eqref{eq:cubical-small-deficit-profile}. For the unit cube, subset weights
are exactly the integers $1,\ldots,2k-1$, which gives
\eqref{eq:cubical-unit-profile} and the stated small-deficit range. If
$M=m_R$, then $\delta=0$ and the upper cap together with
\eqref{eq:cubical-error-zero-mean} forces $q_a=0$ for every face. This gives
the canonical boundary values and completes the proof.
\end{proof}

For $I=\{i_1<\cdots<i_k\}\subset\{1,\dots,n\}$ and
$\varepsilon\in\{0,1\}^{I^c}$, let
\[
\sigma_{I,\varepsilon}
:=
\{x\in I^n:\ x_j=\varepsilon_j \text{ for every } j\notin I\}
\]
be the oriented $k$-cell parallel to
$\dd x_I=\dd x_{i_1}\wedge\cdots\wedge \dd x_{i_k}$, and let
$\delta_{I,\varepsilon}\in C^k(I^n;\RR)$ be the dual basis cochain. Define
\[
\phi_0(t):=1-t,
\qquad
\phi_1(t):=t,
\qquad
\lambda_{I,\varepsilon}(x):=
\prod_{j\notin I}\phi_{\varepsilon_j}(x_j).
\]
The cubical Whitney map is defined by
\[
W(\delta_{I,\varepsilon}):=
\lambda_{I,\varepsilon}(x)\,\dd x_I
\]
and extended linearly. For a smooth $k$-form $\alpha$ on a neighborhood of
$I^n$, the cell-integration map is
\[
\Int(\alpha)(\sigma):=
\int_\sigma \iota_\sigma^\ast\alpha.
\]
Since each $\lambda_{I,\varepsilon}$ is a multilinear polynomial on the whole
cube, $W\theta$ is a globally smooth differential form on $I^n$ for every
cochain $\theta$. Thus the continuous homotopy operator $K_j$ applies without
any regularity issue in the present one-cube setting. The next two
propositions collect the elementary identities needed for the discrete
argument.

\begin{proposition}[Whitney formulas and norm-one bounds]
\label{prop:cubical-whitney}
For the one-cube cubical Whitney map $W$ and cell-integration map $\Int$
defined above,
\[
\dd W\theta=W(\cob\theta)
\qquad
(\theta\in C^k(I^n;\RR)),
\]
and
\[
\Int(\dd\alpha)=\cob(\Int\alpha)
\qquad
(\alpha\in\Omega^r_{\mathrm{sm}}(I^n),\ 0\le r\le n-1),
\]
where both sides are $(r+1)$-cochains. Equivalently, the cell-integration map
commutes with the exterior derivative $\dd$ and the cubical coboundary $\cob$ in
every degree in which both sides are defined. Also,
\[
\Int W\theta=\theta
\qquad
(\theta\in C^k(I^n;\RR)).
\]
Moreover,
\[
\norm{W\theta}_{\coeff,\infty}\le\norm{\theta}_{\cell,\infty},
\qquad
\norm{\Int\alpha}_{\cell,\infty}\le\norm{\alpha}_{\coeff,\infty}
\]
for every cochain $\theta$ and every smooth form $\alpha$ of the corresponding
degree on a neighborhood of $I^n$.
\end{proposition}

\begin{proof}
We first check the Whitney identities on the basis cochains
$\delta_{I,\varepsilon}$. Write
$I=\{i_1<\cdots<i_k\}$ and fix $j\notin I$. Put $J:=I\cup\{j\}$, ordered
increasingly, and let $s(j,J)$ be the position of $j$ in $J$. The oriented
boundary of the cell $\sigma_{J,\eta}$ contains the face
$\sigma_{I,\varepsilon}$ only when $\eta$ agrees with $\varepsilon$ away from
$j$; the incidence sign is
\[
(-1)^{s(j,J)-1}\quad{\rm if}\quad \varepsilon_j=1,
\qquad
-(-1)^{s(j,J)-1}\quad{\rm if}\quad \varepsilon_j=0.
\]
On the other hand,
\[
\dd\bigl(\lambda_{I,\varepsilon}\dd x_I\bigr)
=
\sum_{j\notin I}
\frac{\partial\lambda_{I,\varepsilon}}{\partial x_j}\,
\dd x_j\wedge\dd x_I .
\]
Since
$\partial_{x_j}\phi_1=1$, $\partial_{x_j}\phi_0=-1$, and
$\dd x_j\wedge\dd x_I=(-1)^{s(j,J)-1}\dd x_J$, the coefficient of the
$J$-term is exactly the same incidence sign multiplied by the product of the
remaining inactive factors. This is precisely the coefficient of
$W(\cob\delta_{I,\varepsilon})$, and hence $\dd W=W\cob$.

The identity $\Int\dd=\cob\Int$ is Stokes' theorem with the same
orientation convention. If $\alpha$ has degree $r$ and $\sigma$ is an oriented $(r+1)$-cell, then
\[
(\Int\dd\alpha)(\sigma)
=
\int_\sigma\dd\alpha
=
\int_{\partial\sigma}\alpha
=
(\cob\Int\alpha)(\sigma),
\]
because the cubical coboundary is dual to the oriented cubical boundary,
$(\cob\theta)(\sigma)=\theta(\partial\sigma)$.

For $\Int W=\mathrm{id}$, integrate $W(\delta_{I,\varepsilon})$ over a
$k$-cell $\sigma_{J,\eta}$. If $J\ne I$, the pullback of $\dd x_I$ to
$\sigma_{J,\eta}$ vanishes. If $J=I$, then all variables outside $I$ are fixed
at $\eta$, so
\[
\lambda_{I,\varepsilon}\big|_{\sigma_{I,\eta}}
=
\prod_{j\notin I}\phi_{\varepsilon_j}(\eta_j)
=
\begin{cases}
1,&\eta=\varepsilon,\\
0,&\eta\ne\varepsilon.
\end{cases}
\]
The integral of $\dd x_I$ over the oriented unit cell $\sigma_{I,\eta}$ is
one, so $\Int W$ sends each dual basis cochain to itself.

For the first norm estimate, the coefficient of $\dd x_I$ in $W\theta$ is
\[
\sum_{\varepsilon\in\{0,1\}^{I^c}}
\theta_{I,\varepsilon}\lambda_{I,\varepsilon}(x).
\]
The functions $\lambda_{I,\varepsilon}$ are nonnegative on $I^n$ and their sum
over $\varepsilon$ is
$\prod_{j\notin I}(\phi_0(x_j)+\phi_1(x_j))=1$. Thus each coefficient is a
convex combination of cell coefficients and is bounded by
$\norm{\theta}_{\cell,\infty}$.

Finally, if $\sigma$ is an oriented $k$-cell and $\alpha$ is a smooth $k$-form,
then $\Int\alpha(\sigma)$ is the integral over a unit $k$-cube of one
coordinate coefficient of $\alpha$, up to sign. Its absolute value is therefore
at most $\norm{\alpha}_{\coeff,\infty}$. Taking the maximum over cells proves
the second norm estimate.
\end{proof}

\begin{definition}[Whitney coordinate-monomial cochain]
\label{def:whitney-decomposable}
A cochain $\theta\in C^k(I^n;\RR)$ is \emph{Whitney coordinate-monomial} if
either $\theta=0$, or there is a single ordered coordinate set
$I=\{i_1<\cdots<i_k\}$ such that
\[
\theta=\sum_{\varepsilon\in\{0,1\}^{I^c}}
\theta_{I,\varepsilon}\,\delta_{I,\varepsilon}.
\]
Equivalently, all nonzero $k$-cell coefficients of $\theta$ are supported on
$k$-cells parallel to the same coordinate $k$-plane.
\end{definition}

\begin{proposition}[Characterization of the Whitney coordinate-monomial class]
\label{prop:whitney-characterization}
For $\theta\in C^k(I^n;\RR)$, the form $W\theta$ is coordinate-monomial if
and only if all nonzero $k$-cell coefficients of $\theta$ are supported on
cells parallel to a single coordinate $k$-plane.
\end{proposition}

\begin{remark}
Proposition~\ref{prop:cubical-whitney} gives the Whitney formula and the
norm-one bounds. The support characterization used below identifies exactly
when the sharp continuous estimate applies to $W\theta$: the nonzero
$k$-cell coefficients of $\theta$ must be parallel to one coordinate
$k$-plane.
\end{remark}

\begin{proof}
If $\theta$ is supported on a fixed coordinate set $I$, the formula cited in
Proposition~\ref{prop:cubical-whitney} expresses $W\theta$ as one coefficient
multiplying $\dd x_I$. Hence $W\theta$ is coordinate-monomial.

Conversely, suppose that $\theta$ has a nonzero coefficient
$\theta_{I,\varepsilon}$ and a nonzero coefficient $\theta_{J,\varepsilon'}$
with $I\ne J$. The coefficient of $\dd x_I$ in $W\theta$ is a multilinear
polynomial in the inactive variables for $I$ whose Bernstein coefficient
$\theta_{I,\varepsilon}$ is nonzero. Such a polynomial cannot vanish
identically, since the functions
$\{\lambda_{I,\varepsilon}\}_{\varepsilon\in\{0,1\}^{I^c}}$ are linearly
independent. The same applies to the $\dd x_J$ coefficient. Thus $W\theta$
has at least two nonzero coordinate directions and cannot be
coordinate-monomial. This proves the equivalence.
\end{proof}

\begin{proposition}[One-cube Whitney decomposition]
\label{prop:discrete-stokes-poincare}
Let $1\le k\le n-1$ and let $\theta\in C^k(I^n;\RR)$ be a cochain such that
$\cob\theta$ is Whitney coordinate-monomial. Define
\[
h\theta:=\Int K_k W\theta,
\qquad
\mathcal E\theta:=\Int K_{k+1}W(\cob\theta).
\]
Then
\[
\theta=\cob h\theta+\mathcal E\theta
\]
and
\[
\norm{\mathcal E\theta}_{\cell,\infty}
\le
\frac{1}{2(k+1)}\norm{\cob\theta}_{\cell,\infty}.
\]
The factor $1/(2(k+1))$ is sharp on this one-cube Whitney
coordinate-monomial class.
\end{proposition}

\begin{proof}
By Proposition~\ref{prop:cubical-whitney},
\[
\cob h\theta+\mathcal E\theta
=
\Int\dd K_kW\theta+\Int K_{k+1}W(\cob\theta).
\]
Since $\dd W\theta=W(\cob\theta)$, the homotopy identity gives
\[
\dd K_kW\theta+K_{k+1}W(\cob\theta)=W\theta.
\]
Applying the all-degree identity $\Int\dd=\cob\Int$ from
Proposition~\ref{prop:cubical-whitney} and using $\Int W=\mathrm{id}$ proves the
decomposition.

Because $\cob\theta$ is Whitney coordinate-monomial,
Proposition~\ref{prop:whitney-characterization} shows that $W(\cob\theta)$ is a
coordinate-monomial smooth $(k+1)$-form. Therefore
Theorem~\ref{thm:optimal-continuous-inverse} and
Proposition~\ref{prop:cubical-whitney} give
\[
\norm{\mathcal E\theta}_{\cell,\infty}
\le
\norm{K_{k+1}W(\cob\theta)}_{\coeff,\infty}
\le
\frac{1}{2(k+1)}\norm{W(\cob\theta)}_{\coeff,\infty}
\le
\frac{1}{2(k+1)}\norm{\cob\theta}_{\cell,\infty}.
\]
For sharpness, specialize Theorem~\ref{thm:cubical-boundary-readout} to the
unit cube $I^{k+1}$. Let $\Theta\in C^{k+1}(I^{k+1};\RR)$ be the top cochain
with $\Theta(I^{k+1})=1$, so $\norm{\Theta}_{\cell,\infty}=1$. Choose a
$k$-cochain $\theta$ with $\cob\theta=\Theta$, for instance the canonical
boundary cochain from Theorem~\ref{thm:cubical-boundary-readout}. If the
decomposition theorem held with a constant $C<1/(2(k+1))$, then
\[
\theta=\cob h+E,\qquad \cob E=\cob\theta-\cob^2h=\Theta,\qquad
\norm{E}_{\cell,\infty}\le C.
\]
Thus $E$ would be a cochain primitive of the top cochain with norm below the
sharp lower bound $1/(2(k+1))$ for the unit $(k+1)$-cube and the top cochain
norm $\norm{\Theta}_{\cell,\infty}=1$ in
Theorem~\ref{thm:cubical-boundary-readout}, a contradiction. Therefore the
constant is sharp.
\end{proof}

\begin{remark}
The construction is local to a single cube. On a general cubical complex one
needs additional compatibility arguments across adjacent cubes in order to
produce a global bounded cochain homotopy; compare the algebraic framework in
\cite{Dupont1976,DupontBook1978,Albert2020}.
\end{remark}

\begin{corollary}[Supporting consequence: one-cube one-cochain primitive correction]
\label{cor:hypercube-gradient}
The preceding one-cell decomposition has the following degree-one
primitive-correction form. Let $\theta\in C^1(I^n;\RR)$ and assume that
$\cob\theta$ is Whitney coordinate-monomial. Then
\[
\displaystyle u:=\Int K_1W\theta\in C^0(I^n;\RR)
\]
satisfies
\[
\norm{\theta-\cob u}_{\cell,\infty}
\le
\frac14\norm{\cob\theta}_{\cell,\infty}.
\]
The constant $1/4$ is optimal.
\end{corollary}

\begin{proof}
With the same $u=\Int K_1W\theta$ as in the statement, this is
Proposition~\ref{prop:discrete-stokes-poincare} with $k=1$:
\[
\theta=\cob u+\mathcal E\theta,
\qquad
\norm{\mathcal E\theta}_{\cell,\infty}
\le\frac14\norm{\cob\theta}_{\cell,\infty}.
\]
The optimality is the $k=1$ sharpness statement of that proposition.
\end{proof}


